\label{sec:legal}

\subsection{Physical Geography and Natural Boundaries in Redistricting Law}

The doctrine that redistricting should follow natural geographic boundaries
is among the oldest redistricting criteria in American law.
Before the one-person-one-vote revolution of the 1960s, courts and legislatures
regularly cited rivers, mountain ridges, and coastlines as the appropriate
boundaries for legislative districts.
After \textit{Reynolds v.\ Sims} \citeyearpar{reynolds1964} established
the equal population requirement, natural boundaries could no longer govern
district shape, but they remained a legitimate \emph{secondary} criterion
that plans could follow when doing so was consistent with population equality.

\paragraph{Communities of interest and physical geography.}
In the modern redistricting framework, physical geography arguments are most
commonly presented as a community-of-interest rationale.
Residents sharing a physical landscape share environmental concerns (flood risk,
wildfire exposure, agricultural water rights), infrastructure needs (coastal
access, mountain pass maintenance), and economic dependencies (tourism,
agriculture, fishing) that constitute a genuine community of interest for
legislative purposes.

The NLCD-based M.2 signal operationalizes this doctrine by treating the
physical character of the land cover as a community signal:
tracts that share the same development type (residential, commercial,
agricultural, forested) share a physical community.
Tracts at a development type transition are not a community; they are a
boundary.

\subsection{Water Boundary Doctrine and Case Law}

The hard-cut rule in M.2 (equation~\ref{eq:hardcut}) implements a
well-established legal doctrine: water bodies are natural community
boundaries that should not be crossed by district lines without
geographic justification.

\paragraph{Judicial recognition of water boundaries.}
Several redistricting cases have addressed water boundaries directly.
In \textit{Carstens v.\ Lamm}, 543 F.\ Supp.\ 68 (D.\ Colo.\ 1982),
the district court approved a congressional plan that followed the
Colorado River as a district boundary, noting that the river constituted
a natural geographic division between communities with distinct economic
and environmental interests.
The court observed that plan proponents could rely on natural geographic
features as a neutral, non-partisan rationale for boundary placement.

In \textit{In re Reapportionment of the Colorado General Assembly},
828 P.2d 185 (Colo.\ 1992), the Colorado Supreme Court affirmed that
natural geographic features, including rivers and ridgelines, are among
the legitimate criteria for district boundary placement, and that a plan
that follows such features faces a lower burden of justification than
one that crosses them arbitrarily.

\paragraph{Outer Banks and barrier island communities.}
North Carolina's Outer Banks provides a recurring example in redistricting
litigation and commission proceedings.
The Outer Banks communities---Dare County, Hyde County waterfront, Currituck
Banks---are separated from the mainland by Pamlico Sound and Albemarle Sound,
bodies of water with no fixed bridge crossings for most of their length.
Commission testimony in the 2001, 2011, and 2021 North Carolina redistricting
cycles consistently identified Outer Banks communities as a distinct community
of interest defined by their physical isolation, coastal economy, and
hurricane vulnerability.

Under the M.2 hard-cut rule, Pamlico Sound and Albemarle Sound adjacencies
receive zero edge weight, preventing the \textsc{bisect} algorithm from
assigning mainland tracts and Outer Banks tracts to the same district.
This directly implements the community-of-interest rationale that North
Carolina commissions have accepted in multiple cycles.

\paragraph{Louisiana coastal communities.}
Louisiana's coastal parishes demonstrate the broader application of the
water boundary principle.
The bayou network south of New Orleans creates communities (Isle de Jean
Charles, Cocodrie, Grand Isle) that are physically accessible only by water
or narrow causeways.
Redistricting plans that bundle these communities with inland parishes far
to the north cross physical barriers that define the coastal community's
distinct identity.
The NLCD water classification captures the bayou network's dominant land
cover character: tracts with majority water classification ($p_{\mathrm{wat}} > 0.5$)
along the coastal bayou zone would receive the hard-cut override in any
M.2 application to Louisiana.

\subsection{The Natural Boundary Doctrine}

Beyond water, the broader natural boundary doctrine encompasses any
well-defined physical transition that creates meaningfully distinct
communities on either side.

\paragraph{Mountain ridgelines.}
Colorado's redistricting commission has cited the Continental Divide as a
natural district boundary in every redistricting cycle since 2001.
The Divide separates communities with distinct water rights (Pacific versus
Atlantic watershed), ski industry versus ranching economies, and different
transportation networks.
NLCD captures the mountainous terrain through the forest and shrub/scrub
classifications (codes 41--52) that dominate high-altitude tracts.
A tract that is 90\% evergreen forest ($p_{\mathrm{nat}} \approx 0.90$) has
very low cosine similarity with an adjacent Front Range urban tract
($p_{\mathrm{res}} \approx 0.70$), making the transition a natural
M.2 cut seam.

\paragraph{Agricultural-suburban boundary.}
The transition from cultivated crops (NLCD class~82) and pasture (class~81)
to low-intensity residential development (class~22) marks the rural-suburban
boundary that has been recognized in countless redistricting proceedings as
a meaningful community division.
The California Citizens Redistricting Commission, in its 2011 and 2021
proceedings, consistently drew district boundaries along the
agricultural-suburban edge of the Central Valley, citing the distinct
economic and policy interests of farming communities versus suburban commuter
communities.
The M.2 signal captures this division through the contrast between
$p_{\mathrm{nat}} \approx 0.90$ (agricultural tracts) and
$p_{\mathrm{res}} \approx 0.70$ (suburban tracts).

\subsection{Partisan Neutrality}

NLCD data contains no electoral, racial, or demographic information.
The land cover classification is derived from satellite imagery and
calibrated against ground-truth reference points; it has no knowledge
of the voting history or racial composition of the tracts it classifies.

Any partisan effect of M.2 weights---such as the preference for cutting
at commercial-residential transitions---is a consequence of the geographic
correlation between land use types and partisan voting patterns.
This correlation is an empirical fact about the American population
landscape (urban commercial cores tend to vote differently from
residential suburbs), not a product of the algorithm's choices.
Under \textit{Rucho v.\ Common Cause} \citeyearpar{rucho2019}, the
use of a non-partisan criterion that happens to have partisan effects
does not constitute partisan gerrymandering if the criterion is
independently justified.
Physical land use is an independently justified criterion with deep
roots in redistricting practice.

\subsection{Expert Witness Deployment Template}

Redistricting practitioners deploying M.2 land use weights should
structure expert testimony around the following elements:

\paragraph{1. Data provenance declaration.}
``I used the USGS National Land Cover Database (NLCD) 2021, produced
by the Multi-Resolution Land Characteristics Consortium under contract
with the U.S.\ Geological Survey.
The NLCD classifies the land surface of the contiguous United States
at 30-meter resolution from Landsat satellite imagery.
It contains no voter registration data, election results, racial
composition data, or any other socioeconomic information about the
populations living within the classified areas.''

\paragraph{2. Hard-cut rule justification.}
``The \textsc{bisect} algorithm assigns a zero-weight edge to any pair
of adjacent census tracts where both tracts have majority water
classification (more than 50\% of their land surface classified as
open water by NLCD).
This prevents the algorithm from placing such tracts in the same district.
This encoding reflects the longstanding redistricting principle that
water bodies are natural community boundaries.
The water classification is objective and derived from satellite imagery,
not from any political or demographic data.''

\paragraph{3. Community boundary identification.}
``In [state], the following physical land use transitions are preserved as
district boundaries by the proposed plan: [specific examples, e.g.,
`the Pamlico Sound water boundary separates the Outer Banks tracts in
Dare County from the mainland, consistent with the commission's recognition
of the Outer Banks as a distinct coastal community in its 2011 and 2021
redistricting cycles'].''

\paragraph{4. Criterion alignment.}
``The [state]'s redistricting criteria require preservation of communities
defined by geographic and natural boundary considerations.
The land use character weights directly implement this criterion: district
cuts run preferentially at land use transitions that correspond to recognized
natural boundaries between communities.''

\subsection{When Not to Use M.2}

Land use weights should not be used in the following contexts:

\begin{itemize}
  \item \textbf{Dense urban cores.}
    In highly urbanized census tracts where most of the land is class~24
    (high-intensity developed), the M.2 vector becomes near-degenerate
    ($p_{\mathrm{com}} \approx 1.0$ for all tracts).
    All tracts look similar under M.2, and the signal adds no information.
    M.1 (economic character) is more informative in this context.
  \item \textbf{States with no significant natural boundaries.}
    In states with uniform terrain and no significant water bodies
    (e.g., Kansas, Iowa), the M.2 signal will produce modest or zero
    effect.
    The agricultural uniformity of the Great Plains means most adjacent
    tracts have similar NLCD vectors; there are no sharp land use
    transitions to exploit.
  \item \textbf{When hard cuts conflict with VRA compliance.}
    In rare cases, the hard-cut rule may prevent contiguous assignment
    of tracts needed for a majority-minority district.
    In such cases, the hard-cut rule must be overridden via bridge
    adjacency encoding, and the VRA requirement takes priority.
\end{itemize}
